\chapter{Building blocks: add, subtract, multiply by one limb}
\label{ch:basics}

Every big algorithm in this library eventually does the same few things over and over:
add two arrays of limbs, subtract one from another, and add a multiple of one array to
another. These operations are called \emph{kernels}. They live in
\file{include/biginteger/common/LimbKernels.h}. This chapter shows how they work and
why the way they are written matters so much.

\section{Addition with carries}

Adding big numbers is grade-school addition. Add the lowest digits, write down the
result, carry one if it overflowed, move left.

\begin{example}[: the same idea in base 10 and base $2^{64}$]
Base 10, $478 + 256$:
\[
\begin{array}{r@{\quad}l}
8 + 6 = 14 & \text{write } 4, \text{ carry } 1\\
7 + 5 + 1 = 13 & \text{write } 3, \text{ carry } 1\\
4 + 2 + 1 = 7 & \text{write } 7
\end{array}
\qquad\Rightarrow\qquad 734
\]
Base $2^{64}$, adding $[\,2^{64}-1,\ 1\,]$ and $[\,1\,]$ (that is, $(2^{64}+2^{64}-1) + 1$):
\[
\begin{array}{r@{\quad}l}
(2^{64}-1) + 1 = 2^{64} & \text{write } 0, \text{ carry } 1\\
1 + 0 + 1 = 2 & \text{write } 2
\end{array}
\qquad\Rightarrow\qquad [\,0,\ 2\,] = 2 \cdot 2^{64}.
\]
\end{example}

A carry is never more than 1, because the largest possible sum,
$(B-1) + (B-1) + 1 = 2B - 1$, is less than $2B$. Subtraction works the same way with a
\emph{borrow} of 0 or 1.

\section{Multiply an array by one limb}

The workhorse of multiplication and division is ``add $b$ times the array $a$ into the
array $r$'', written \code{AddMul1(r, a, n, b)}. Each step multiplies one limb $a_j$ by
$b$, which gives a two-limb (128-bit) product: a low half and a high half. The low half
goes into position $j$; the high half spills into position $j+1$.

\begin{example}[: $987 \times 6$, digit by digit]
Little-endian digits of 987 are $[7, 8, 9]$. Multiply each by 6:
\[
7\cdot 6 = 42,\qquad 8 \cdot 6 = 48,\qquad 9 \cdot 6 = 54 .
\]
Each product has a low digit (2, 8, 4) and a high digit (4, 4, 5). Position $k$ of the
answer is low digit $k$ plus high digit $k-1$ plus any carry:
\[
\begin{array}{c|cccc}
\text{position} & 0 & 1 & 2 & 3\\ \hline
\text{low of } p_k & 2 & 8 & 4 & \\
\text{high of } p_{k-1} & 0 & 4 & 4 & 5\\
\text{carry in} & 0 & 0 & 1 & 0\\ \hline
\text{digit} & 2 & 2 & 9 & 5
\end{array}
\]
Result $[2,2,9,5] = 5922 = 987 \times 6$. In the library the ``digits'' are 64-bit limbs
and the products are 128-bit, but the bookkeeping is identical.
\end{example}

\section{Why the way you write the loop matters}
\label{sec:flagchain}

The obvious C++ for \code{AddMul1} is

\begin{lstlisting}
unsigned __int128 p = (unsigned __int128)a[j] * b + r[j] + carry;
r[j]  = (uint64_t)p;          // low 64 bits
carry = (uint64_t)(p >> 64);  // high 64 bits feed the next limb
\end{lstlisting}

It is correct, and it was in the library for years. The trouble is the
\code{carry} variable. It holds a whole limb, and the next step cannot begin adding until
it is known. Every step waits for the previous step's 128-bit arithmetic to finish. A
modern processor can start a new multiplication every cycle, but this loop forces it to
work one limb at a time, at 2.5--3.3 cycles per limb.

The kernels in \file{LimbKernels.h} split the work into \textbf{two separate carry
chains}, as in the worked example above:

\begin{enumerate}
  \item \textbf{Chain A} forms the product row: $t_k = \text{low}(p_k) + \text{high}(p_{k-1})$,
        with a one-bit carry.
  \item \textbf{Chain B} adds that row into the result: $r_k \mathrel{+}= t_k$, with its own
        one-bit carry.
\end{enumerate}

Neither chain waits on a full 128-bit value. Both carries are single bits, and the code is
written in straight-line blocks of four limbs with the compiler's carry builtins
(\code{\_\_builtin\_addcll}), so the carry bit stays in the processor's flags register from
one limb to the next. On ARM this compiles to a run of \code{adcs} (add with carry)
instructions, the same instructions a hand-written assembly routine would use.

\begin{figure}[ht]
\centering
\begin{tikzpicture}[x=1.15cm,y=1cm]
  \node[lbl,anchor=east] at (-0.5,1.6) {one chain};
  \foreach \k in {0,...,5} {
    \node[box,fill=sBlue!12,minimum width=11mm] (o\k) at (\k*1.3,1.6) {\footnotesize mul+add};
  }
  \foreach \k [evaluate=\k as \n using int(\k+1)] in {0,...,4} \draw[arr,sRed] (o\k) -- (o\n);
  \node[lbl,anchor=west] at (7.3,1.6) {each step waits for a 64-bit carry};
  \node[lbl,anchor=east] at (-0.5,0.4) {chain A};
  \node[lbl,anchor=east] at (-0.5,-0.6) {chain B};
  \foreach \k in {0,...,5} {
    \node[box,fill=sOrange!15,minimum width=11mm] (a\k) at (\k*1.3,0.4) {\footnotesize lo+hi};
    \node[box,fill=sAqua!15,minimum width=11mm] (b\k) at (\k*1.3,-0.6) {\footnotesize $r \mathrel{+}= t$};
    \draw[arr,muted] (a\k) -- (b\k);
  }
  \foreach \k [evaluate=\k as \n using int(\k+1)] in {0,...,4} {
    \draw[-{Stealth[length=1.5mm]},sGreen,thick] (a\k) -- (a\n);
    \draw[-{Stealth[length=1.5mm]},sGreen,thick] (b\k) -- (b\n);
  }
  \node[lbl,anchor=west,align=left] at (7.3,-0.1) {two chains of 1-bit carries\\run side by side};
\end{tikzpicture}
\caption{Dependency chains in \code{AddMul1}. Red: the old loop passes a full limb from
step to step. Green: the new kernel passes only carry bits, along two independent
chains that the processor overlaps.}
\end{figure}

\begin{perfnote}
Cycles per limb on an M1 core:
\begin{center}
\begin{tabular}{lccc}
\toprule
kernel & old loop & GMP (hand assembly) & new kernel \\
\midrule
\code{addmul\_1} & 2.5--3.3 & 1.3--1.8 & 1.64--1.76 \\
\code{add\_n} & 2.2--3.3 & 0.65--0.9 & 1.1--1.7 \\
\code{submul\_1} (division) & 2.9--4.0 & 1.3--1.8 & 1.6--1.8 \\
\bottomrule
\end{tabular}
\end{center}
Schoolbook multiplication became 1.5--1.8 times faster, with no assembly language
anywhere in the library. Earlier the gap to GMP had been blamed on ``portable C++ cannot
match assembly''. It turned out the gap was mostly how the loop was written.
\end{perfnote}

\begin{pitfall}
Before this change, an attempt to speed up schoolbook by processing two rows at once
(``2-row unrolling'') measured no gain at all. It kept the same limb-valued carry, so
it kept the same bottleneck. Speed came from changing the \emph{dependency structure}, not
from doing more work per loop iteration.
\end{pitfall}

\section{The kernel set}

\begin{center}
\small
\begin{tabularx}{\textwidth}{lX}
\toprule
kernel & what it computes \\
\midrule
\code{AddN}, \code{SubN} & $r = a \pm b$, arrays of equal length, returns carry / borrow \\
\code{Mul1} & $r = a \cdot b$ for a single limb $b$ \\
\code{AddMul1} & $r \mathrel{+}= a \cdot b$ (one row of schoolbook multiplication) \\
\code{SubMul1} & $r \mathrel{-}= a \cdot b$ (one step of long division) \\
\code{MulBasecase} & schoolbook product: one \code{Mul1} row then \code{AddMul1} rows \\
\code{SqrBasecase} & schoolbook square (\cref{sec:square}) \\
\code{Lshift}, \code{Rshift} & shift an array left or right by 0--63 bits \\
\bottomrule
\end{tabularx}
\end{center}

These kernels are used only with true 64-bit limbs. The legacy layout
(\code{-DBIGMATH\_LIMB\_64=0}) keeps its own simpler loops.

\newpage
\section{Parallel addition and subtraction}
\label{sec:paralleladdsub}

A carry or borrow can cross the entire number. BigMath splits the magnitude into
contiguous blocks and computes each block with an incoming flag of zero. It then
uses the provisional results to determine which blocks need correction. The sign
handling still happens in the operator layer, so adding opposite signs can use
the magnitude subtraction kernel.

\begin{enumerate}
  \item Each block runs the ordinary limb kernel and saves its outgoing carry
        or borrow. For addition it finds the run of all-ones limbs at the low end
        of its provisional result. For subtraction it finds the run of zeros.
  \item A short serial pass visits blocks from low to high. Each block passes
        along its own outgoing flag, or an incoming flag if the run covers the
        entire block. This determines the incoming flag for every block.
  \item A block receiving a carry turns its run of all-ones limbs into zeros
        and increments the next limb, if present. A block receiving a borrow
        turns its zeros into all-ones limbs and decrements the next limb. If at
        least two corrections have runs of 16\,384 limbs or more, the pool
        splits the fixup work again. Otherwise the caller applies the corrections.
\end{enumerate}

\begin{example}
Split $9999 + 1$ into two-digit blocks. The low block computes $99+1$, writes
$00$ and sends a carry. The high block provisionally writes $99$. Its incoming
carry turns both digits to zero and continues beyond the high block, giving $10000$.
\end{example}

In a threaded 64-bit build, the dispatcher considers this path for equal-length
magnitudes above the stack-buffer path, with at least 2\,048 limbs per block.
Addition and subtraction each have a warm and a parked cutoff, measured in
limbs per operand. The checked-in profiles inherit zero for all four cutoffs,
so their magnitude addition and subtraction run serially. The kernel is in
\file{include/biginteger/algorithms/ParallelAddSub.h}.

\subsection*{Choosing the cutoffs}

The tuner requires a speedup of at least $1.05\times$ on both random operands and
full carry or borrow chains at two adjacent sizes and every larger sampled size.
The workflow selects the highest cutoff supported by all three surveys. The hosted
survey selected zero on all five release targets. A native M1 Max survey gave
these machine-specific cutoffs, in limbs per operand.

\begin{center}
\small
\begin{tabular}{lrr}
\toprule
Magnitude operation & Warm & Parked \\
\midrule
Addition    & 24\,576 & 98\,304 \\
Subtraction & 16\,384 & 1\,048\,576 \\
\bottomrule
\end{tabular}
\end{center}

In the native holdout on random equal-length inputs, the public operators' geometric
mean speedups for addition and subtraction were $1.77\times$ and $2.24\times$ in
bursts and $1.54\times$ and $1.09\times$ with parked workers. The parked borrow-chain
group averaged $0.97\times$. See \file{BENCHMARK.md}, \file{docs/PLATFORM\_TUNING.md}
and the build macros in \cref{app:ref}.

\newpage
\section{Comparison and shifts}

For numeric ordering, \code{CompareTo} checks the signs first and reverses
the magnitude order when both values are negative. The relational operators
\code{>}, \code{>=}, \code{<}, and \code{<=} use this same result.

The raw \code{Compare} routine treats high zero limbs as padding. If the
stored lengths differ, any nonzero limb in the longer operand's excess
range decides the comparison. Otherwise, compare the common width. Check
its two highest limbs individually, then scan the remaining limbs in
blocks of 16 when the range is long enough. In each block, XOR corresponding
limbs and OR the differences together. A zero result skips the entire block.
In the first differing block, compare whole limbs from most to least
significant. The same block scan can recognize shared storage, and raw
comparison with a scalar can skip blocks of high zeros.

\begin{example}
Suppose the highest difference is $0\mathrm{x}0100$ on the left and
$0\mathrm{x}00\mathrm{ff}$ on the right, and a lower limb is smaller on
the left. The left value is larger because the highest differing limb
decides. Comparing the first unequal bytes in memory would give a
different result on a little-endian machine.
\end{example}

The worst case remains $O(n)$ time and $O(1)$ extra space.
In the M1 Max benchmark, an ordered comparison of 4,096 limbs with a
low-limb mismatch took about 0.45\,\textmu s, down from
1.30\,\textmu s. Comparisons settled by different lengths took
about $0.7$\,ns longer. See \file{BENCHMARK.md} for the cases and commands.

For \code{a == b}, comparing an object with itself takes constant time. For distinct
objects, the operator checks the signs, stored lengths and outermost limbs. If lengths
differ and the longer value has a high zero limb, it compares the magnitudes to account
for padding. For long values of the same width, it checks the lowest eight limbs as a
block, then compares eight limbs at a time from the high end. This catches small
low-order changes immediately and still exits early for differences near the top.
Equal copies require an exact scan of all $n$ limbs, taking $O(n)$ time and $O(1)$
extra space. The \code{!=} operator uses the same equality check.

The public shift operators count whole limbs. \code{a << k} multiplies
by $B^k$, and \code{a >> k} drops the lowest $k$ limbs while preserving
the sign of a nonzero result. Internally, a shift by $s$ bits splits into
$\floor{s/64}$ whole limbs and $s\bmod 64$ bits in the default radix.
Division uses those bit shifts to normalize the divisor
(\cref{sec:normalize}).

The shared limb utilities use \code{DataT} for stored limb values and
\code{SizeT} for the library's limb counts. Container offsets use
\code{std::size\_t}. Division shares a leading-zero count for
normalization and a trimmed limb count for padded operands.
